\section{Capsule three: subsequences, and a plan for the proof}\label{ch17:sec:subsequence}
\par\noindent\begin{minipage}{\linewidth}
\begin{definition}[Subsequence]\label{ch17:def:subsequence}
Let $x_0,x_1,x_2,\ldots$ be a sequence. A \emph{subsequence} of it is a sequence of the form
\[
 x_{n_1},x_{n_2},x_{n_3},\ldots,
 \qquad\text{where }n_1<n_2<n_3<\cdots
\]
are natural numbers.
\end{definition}
\end{minipage}\par

In words: we may skip terms, but the terms we keep must stay in their original order, and no index may be used twice. For $x_n=(-1)^n$, the indices $n_k=2k$ give the subsequence $x_2,x_4,x_6,\ldots$, which is the constant sequence $1,1,1,\ldots$. By contrast, $x_1,x_0,x_3,x_2,\ldots$ is not a subsequence, because its indices do not increase, and neither is $x_0,x_0,x_0,\ldots$, which uses the index $0$ again and again.

Keep the two counters apart: $k$ is the position in the subsequence, while $n_k$ is the position of the same term in the original sequence; in the example, the third term of the subsequence is $x_6$. We number the terms of a subsequence from $k=1$, not $0$. This makes no difference: with $y_k=x_{n_k}$, the definitions of convergence and of Cauchy sequences apply to $y_1,y_2,y_3,\ldots$ as before, with thresholds that are at least $1$.

To make the chosen indices increase, we need one fact. \emph{An infinite set $S$ of natural numbers contains a number larger than any given natural number $M$.} Otherwise every member of $S$ would lie in $\set{0,1,\ldots,M}$, and crossing out of the list $0,1,\ldots,M$ every number not in $S$ would leave a list of the members of $S$ that ends, so $S$ would be finite. So if each stage of a construction offers an infinite set of allowed indices, we can always choose an allowed index larger than the one chosen at the previous stage.

\subsection*{A plan with three kinds of objects}
The construction involves three kinds of objects that must not be mixed up: points of the space, such as the term $x_7$; balls, which are sets of points; and sets of indices, which are sets of natural numbers. If $x_7=x_9$ lies in a chosen ball, it is one point of the ball but contributes two indices, $7$ and $9$.

Here is the plan. Start with every index allowed: $S_0=\NN$. At stage~$1$, cover the space by finitely many balls of radius $1/4$, as total boundedness allows. By Lemma~\ref{ch17:lem:pigeonhole}, one of these balls contains $x_n$ for infinitely many indices $n\in S_0$; keep exactly those indices, as the set $S_1$. By the diameter bound, any two terms with indices in $S_1$ are less than $2\cdot\tfrac14=\tfrac12$ apart. In general, stage~$k$ uses balls of radius $2^{-(k+1)}$, considers only the indices kept at stage~$k-1$, and keeps infinitely many indices whose terms are less than $2\cdot2^{-(k+1)}=2^{-k}$ apart. The radius is half the desired distance because the diameter bound doubles it; any radii shrinking towards zero would work, and this choice keeps the arithmetic tidy. Since each stage only throws indices away,
\[
 S_0\supseteq S_1\supseteq S_2\supseteq\cdots,
\]
where $A\supseteq B$, read ``$A$ contains $B$,'' means $B\subseteq A$. We say that the index sets are \emph{nested}. The chosen balls need not be nested: the ball of stage~$2$ may stick out of the ball of stage~$1$. That does no harm, because a term whose index lies in $S_2$ also has its index in $S_1$, so it lies in both chosen balls. The nesting of the index sets carries the distance bound of each stage forward to all later stages.
